\documentclass[12pt]{article}
\usepackage[utf8]{inputenc}
\usepackage[T1]{fontenc}
\usepackage{hyperref}
\title{Sos Kundi Language Resources for K-12 Curriculum Development}
\author{Compiled by Claude Sonnet 5.5 (claude-sonnet-5-5) for Andr\'as Kornai}
\date{}
\begin{document}
\maketitle

\section{Introduction}
Sos Kundi (ISO 639-3 \texttt{sdk}; Glottocode sosk1235) is a Ndu language of East Sepik Province, Papua New Guinea. Sources found in search give about 3,500 speakers (2018 census, per Wikipedia), and other sources give different figures. The Ndu family is described in the Wikipedia article on the Ndu languages. This dossier is short because little published material exists beyond a dictionary and some audio and Bible portions. Evidence of use in schools was not found.

\section{Primary Educational Resources}
No primer, textbook or school curriculum for Sos Kundi turned up in the searches. The closest items are the dictionaries listed below, which could serve as the base for early-literacy materials.

\section{Linguistic and Reference Materials}
\begin{itemize}
\item Sos Kundi dictionary on Webonary, compiled by Correna Janzen and Elvira Corbalan (SIL) with community members over about ten years. The compilers call it a first draft. It is searchable in Sos Kundi, Tok Pisin and English. \url{https://www.webonary.org/soskundi/}
\item Sos Kundi Trilingual Dictionary (Sos Kundi--English--Tok Pisin), PDF, licensed CC BY-NC-SA 4.0. \url{https://www.webonary.work/soskundi/files/Sos-Kundi-Trilingual-Dictionary.pdf}
\item Wikipedia, Sos Kundi language. \url{https://en.wikipedia.org/wiki/Sos_Kundi_language}
\item Wikipedia, Ndu languages (family context). \url{https://en.wikipedia.org/wiki/Ndu_languages}
\item ASJP word-list entry for Sos Kundi. \url{https://asjp.clld.org/languages/SOS_KUNDI}
\end{itemize}

\section{Religious and Community Resources}
Joshua Project reports Bible portions translated in 2014 and no complete New Testament or Bible. Its figures for the population differ between versions of the profile. \url{https://joshuaproject.net/people_groups/14861}

\section{Digital and Audio Resources}
Global Recordings Network lists Sos Kundi recordings. \url{https://globalrecordings.net/en/language/sdk}

\section{Curriculum Development Context}
The searches found no information on Sos Kundi teaching in schools. Any curriculum would have to start from the SIL dictionary and from collaboration with the community and SIL. The Tok Pisin glosses in the trilingual dictionary suit the PNG vernacular-then-Tok-Pisin-and-English approach to schooling. The dictionary is marked as a draft, so its orthography and entries would need checking by speakers.

\end{document}